\chapter{Reproducibility Record}\label{app:repro}
The code is in the public repository \url{https://github.com/mrchowdhury587/foodprice_watch}. Each notebook clones the code, runs on the data folder, and prints its code commit and the sha256 hashes of its inputs and outputs. Notebooks 02 to 06 were run on the work computer and again in Google Colab; every output was byte-identical (Table~\ref{tab:repro}).

\begin{table}[H]\centering\small
\caption[Output files and their hashes]{Main output files and the first eight characters of their sha256 hash, identical in the local and the Colab runs.}\label{tab:repro}
\fittab{\begin{tabular}{lll}
\toprule
\textbf{File} & \textbf{Written by} & \textbf{sha256 (first 8)} \\
\midrule
\texttt{data\_8.csv} & 02 processing & 381ef71d \\
\texttt{rise\_episodes\_8.csv} & 03 price behaviour & 48d529fc \\
\texttt{forecasts\_test.csv} & 04b forecast test & d536b503 \\
\texttt{forecasts\_h7.csv} & 04b forecast test & 6449ec4d \\
\texttt{alerts\_8.csv} & 05b alert scoring & 613b3c45 \\
\texttt{R1\_data.csv} & 06 report tables & 63cafc13 \\
\bottomrule
\end{tabular}
}
\source{\path{reports/RUN_LOG.md}.}
\end{table}

\begin{table}[H]\centering\small
\caption[Commits that fixed the settings]{Commits that fixed the settings before scoring.}\label{tab:commits}
\begin{adjustbox}{max width=\linewidth}\begin{tabular}{lll}
\toprule
\textbf{Commit} & \textbf{What was fixed} & \textbf{Before} \\
\midrule
\texttt{8d7c145} & LightGBM: price only, 7 leaves, rate 0.03, absolute-error loss & notebook 04b (forecast test) \\
\texttt{e0cfa06} & alert settings: ML range 0.95, simple range 0.95, fixed rise 7\,\% & notebook 05b (alert scoring) \\
\bottomrule
\end{tabular}\end{adjustbox}
\end{table}

\chapter{Results Register}\label{app:register}
Every headline number of this report is listed in \path{reports/RESULTS_REGISTER.md} with the committed table it comes from, the notebook that wrote the table, and the commit. Results for 2025-26 carry the label ``retrospective''.

\chapter{Changes to the Original Plan}\label{app:deviations}
Every change to items, dates, drivers, method, thresholds or sources was written into \path{DEVIATIONS.md} before work continued. The list below covers every entry; the file gives the reasons and approvals.
\begin{itemize}[leftmargin=*,itemsep=2pt]
\item \textbf{Scope.} The project was restarted under the present title. The earlier pipeline stays unchanged in the archive. The forecasts became a benchmark and the monitoring and alert system became the product, because the earlier model had not beaten \naive{}.
\item \textbf{Items.} Eight items instead of 54, chosen for how often their prices move and for being everyday purchases, after the earlier results had been seen (Section~\ref{sec:threats}).
\item \textbf{Drivers.} Kept: calendar, exchange rate and weather. Dropped: the India onion export flag (its source ends in March 2025), public holidays (the Eids are in the Ramadan--Eid file) and news events.
\item \textbf{Exchange rate.} One source only (Yahoo), without the earlier splice. The 30-day change now takes the last quote on or before the day 30 days earlier (at most 7 days older); looking up that exact day had left 28\,\% of values empty. The Yahoo series is stale from about September 2026, so the 30-day change is understated in the last weeks.
\item \textbf{Data files.} Four raw files; the item list moved into the settings file; the panel and feature files were merged into one table of 20 columns with no change of values.
\item \textbf{Forecast.} Three candidates (\naive{}, ETS, LightGBM) instead of six; LightGBM chosen from a 16-configuration grid on 2024 instead of a 50-trial search; 82 weekly test origins whose 30-day target lies on or before 24 October 2026.
\item \textbf{Alerts.} Each rule's setting is the most sensitive candidate whose 2024 alert runs that match no sustained rise stay at or below 2 per item-year; the earlier run counted alerts outside news-event windows, which were dropped.
\item \textbf{Test period.} 2025-26 had been scored by the earlier run, so every result on it is labelled retrospective; all choices were still made on 2024.
\item \textbf{Plan claims not met.} Alerts catch 55 to 74\,\% of rises (plan: 65 to 75\,\%), 2.5 to 2.8 times the random-date rate (plan: about three), with three forecast models (plan: six). Two-sided range coverage is not reported because the ranges are upper bounds only.
\item \textbf{Runs.} Notebooks 03 to 06 were first run on the work computer while Google Drive could not be mounted in Colab, then rerun in Colab with identical outputs; the Colab runs are the runs of record.
\item \textbf{Monitor.} Notebook 07 and the Gradio app replay the three rules to 24 October 2026. The app can be shared through a temporary public link, approved by the project owner.
\item \textbf{Report.} This report replaces the earlier one on the same university template: new title and four research questions; the driver-breakdown and news-event chapters were dropped.
\end{itemize}
